\section{第\ref{sec:livingmachine}章　生きたニューロンでできたマシン}

電極アレイ上の培養の振る舞いは、記録と刺激を伴う直接の実験で測定されている。一斉発火の機構はシナプス枯渇のモデルと微小培養での実験に基づき、シャーレの中のドーパミンに関する主張は単発の実験に基づく。そのうちいくつかは著者ら自身がデモンストレーションにすぎないとみなしている。

{\footnotesize
\setlength{\tabcolsep}{3pt}\renewcommand{\arraystretch}{1.15}
\begin{longtable}{>{\raggedright}p{0.5cm}>{\raggedright}p{4.0cm}>{\raggedright}p{5.6cm}>{\raggedright}p{2.3cm}>{\raggedright\arraybackslash}p{1.7cm}}
\toprule
No. & 主張 & 実験と測定されたもの & 出典 & 状態 \\
\midrule\endfirsthead
\toprule
No. & 主張 & 実験と測定されたもの & 出典 & 状態 \\
\midrule\endhead
1 & チップ上の培養：大部分は\nt{GLUT}、少数の\nt{GABA}の割合は標本による。海馬の培養では約 $6\%$ & 電極アレイ上の数千ニューロンのネットワークは標準的な実験である（行~3）。海馬の培養での\nt{GABA}ニューロンの割合は GABA 染色で数えられ、ニューロンの約 $6\%$ だった。大脳皮質の培養については検証された数値を挙げない & \cite{Benson1994} & 測定済み \\
2 & オルガノイド：ヒトの幹細胞からできた、大きさ約1ミリの三次元神経組織 & 幹細胞から、複数の脳領域をもつ三次元オルガノイドが育てられた。その中には前駆細胞の層と成熟ニューロンをもつ大脳皮質も含まれる & \cite{Lancaster2013} & 測定済み \\
3 & 入力がないと培養は、培養によって1秒から数分の間隔で一斉発火する & 密度 $3\,000$ から $50\,000$ ニューロンの大脳皮質培養 $58$ 個を5週間記録した（30分の記録 $963$ 本）。2週間後にはほぼすべての培養でネットワーク全体の一斉発火が優勢になる。一斉発火の間隔は $1$ から $300$~s で、標本に強く依存する & \cite{Wagenaar2006} & 測定済み \\
4 & 一斉発火は伝達物質の枯渇で終わり、間隔は $T_{\text{burst}}\approx\tau_{\text{rec}}\ln\frac{1}{1-x^*}$~\eqref{eq:burstinterval}。$2$--$4$~s は $\tau_{\text{rec}}=0.8$~s のモデルによる推定で、測定された回復はもっと遅い & 式は本章で導いた。モデル：枯渇するシナプスをもつネットワークは、自らネットワーク全体の一斉発火を生む。$4$--$30$ ニューロンの微小培養での実験：一斉発火の持続時間は前回からの時間に依存し、伝達物質の小胞を枯渇させると一斉発火は消える。著者らの結論は、一斉発火は伝達物質の蓄えで制限されるというもの。そこでの蓄えの回復は数十秒で、同じ式で数十秒から数分の間隔を与える & 本章の導出、\cite{Tsodyks2000,CohenSegal2011,Tsodyks1997} & 理論 \\
5 & 「黙る教師」：目的の応答が出たら刺激を止め、応答が定着する & 培養を $0.3$--$1$~Hz で刺激し、刺激後 $50\pm10$~ms に選んだ応答が現れたら刺激を $5$~min 止めた。数サイクル後、目的の応答はすぐに誘発されるようになった。学習曲線が描かれた & \cite{Shahaf2001} & 測定済み \\
6 & 培養がテニスをする。セッション内のラリーの有意な伸びは完全なフィードバックのときだけ & 詳細は第\ref{sec:dishbrain}章とその補遺。刺激の代わりに無音にしても、セッションの終わりにはフィードバックなしより上手になるが、効果は小さい。セッション間の学習は不安定 & \cite{Kagan2022} & 測定済み \\
7 & 培養が自分の入力を予測するように学習するというのは仮説。「知性」をめぐる大げさな言葉には反論がある & 自由エネルギー原理は理論的な提案であり、より単純な説明と区別する直接の実験はない。30人の研究者が反論論文で、ループ内での振る舞いの変化は知性も感覚も示さないと指摘した & \cite{Friston2010,Balci2023} & 仮説 \\
8 & 刺激パターンを選べば、培養は仮想の体を目標へ動かす & プログラムが現在の成績に応じて訓練刺激のパターンを自動で選んだ。数十分でネットワークは指定した状態に達し、$2$~h の訓練後、可塑性は $80$~min 続いた。振る舞いと無関係に与えた同じ刺激は効果がなかった & \cite{Bakkum2008} & 測定済み \\
9 & リザバー：学習するのは線形読み出し $y=W_{\text{out}}\mathbf r$~\eqref{eq:reservoir} だけ & リザバー計算の理論：減衰する記憶をもつ任意の非線形システムに、学習させた線形出力を加える & \cite{Maass2002,JaegerHaas2004} & 理論 \\
10 & リザバーとしてのオルガノイドは話者を約 $78\%$ の精度で区別する（著者らのデータ）。誤差で学習するのは読み出しで、オルガノイドの結合は教師なしで変わる & 複数の話者の音声を、電極アレイ上の電流パターンとしてオルガノイドに与えた。学習させた読み出しは約 $78\%$ の精度で話者を区別した。これは著者らの報告である。著者らはまた、訓練中にオルガノイドの機能的結合が変化したと報告し、これをオルガノイド内の教師なし学習と呼んでいる & \cite{Cai2023} & 測定済み \\
11 & 100万個のニューロンがスパイクに使う電力は約 $10^{-4}$~W~\eqref{eq:dishpower} & 推定：大脳皮質のエネルギー収支\eqref{eq:energy}から得たスパイクのコスト $10^{-10}$~J にスパイク数を掛けた。培養の電力の直接測定はない & 本章の導出、\cite{Attwell2001} & 理論 \\
12 & 光のフラッシュで放出するドーパミン：ケージドドーパミン $1$~mM、$240$ 回の繰り返しで誘発スパイク数が増加 & オルガノイドの遠隔プラットフォームで、刺激がより多くのスパイクを誘発したとき、光感受性ケージに入ったドーパミン（$1$~mM）を紫外線（$365$~nm、$0.8$~s）で解放した。$240$ ステップ（約 $5$~h）でスパイク数は全体として増えた。著者らは、実験一つではドーパミンとの関係は確立できないと書いている。対照はない & \cite{Jordan2024} & 仮説 \\
13 & 生きた細胞からのドーパミンが有機トランジスタの重みを長期かつ可逆的に変える & ドーパミンを分泌する細胞を有機ニューロモルフィック素子の上で育てた。電極でのドーパミンの酸化がチャネルのコンダクタンスを変える。重みの長期変化とその復帰が示された & \cite{Keene2020} & 測定済み \\
14 & 働く\nt{DOPA}ニューロンをもつオルガノイドは存在する。そのドーパミンを教師として使った例はない & ヒトの幹細胞からのオルガノイドに、電気的に活動する成熟\nt{DOPA}ニューロンとドーパミン産生が見つかった。このドーパミンが別のネットワークを教える実験は文献に見つからなかった & \cite{Jo2016} & 測定済み \\
15 & オルガノイドが棒を立てる。プログラムによる学習刺激はランダムより良いが定着せず、\nt{GLUT}シナプスを介する & マウス大脳皮質のオルガノイドを「台車の上の棒」課題と閉ループでつないだ。大多数のオルガノイドで、強化学習が選んだ信号はランダムな信号や信号なしより良い成績をもたらした。$45$~min の休止後、改善は残らなかった。AMPA 受容体と NMDA 受容体の遮断で改善は消えた & \cite{Robbins2026} & 測定済み \\
16 & 制御レジスタがなければ、テストベンチ上のネットワークは何を成功とみなすかを自分で決められない（命題\ref{prop:livingmachine}） & 行 5--15 のすべての実験で、学習信号は実験者かプログラムが与える。培養が自ら成功の基準を作り出した実験はない。これは不在からの推論であり、実験ではない & --- & 仮説 \\
\bottomrule
\end{longtable}}
